\documentclass[10pt]{article}
\usepackage[margin=1in]{geometry}
\usepackage{amsmath,amssymb,amsthm,booktabs,microtype}
\usepackage[hidelinks]{hyperref}
\linespread{1.05}

\theoremstyle{plain}
\newtheorem{theorem}{Theorem}
\newtheorem{proposition}[theorem]{Proposition}
\newtheorem{corollary}[theorem]{Corollary}
\theoremstyle{definition}
\newtheorem{definition}[theorem]{Definition}
\newtheorem{assumption}[theorem]{Assumption}
\newtheorem{remark}[theorem]{Remark}

\newcommand{\R}{\mathbb{R}}
\newcommand{\Ex}{\mathbb{E}}

\title{\vspace{-2em}What Content-Adaptive Early Exit Is Worth:\\
An Exact Characterisation of the Compute--Distortion Frontier}
\author{}
\date{}

\begin{document}
\maketitle
\vspace{-3em}

\begin{abstract}
\noindent
A multi-exit neural decoder can spend different amounts of computation on
different regions of one frame. We give an exact description of what such a
decoder can achieve. Writing $C_k$ for the relative cost of exit $k$ and
$D_i(k)$ for the distortion tile $i$ incurs there, we show that the achievable
(cost, distortion) set under \emph{fixed-proportion mixing} is precisely the
convex hull of the $K$ per-exit points (Proposition~\ref{prop:hull}), while
under \emph{per-tile assignment} it is the Minkowski average of the $N$
per-tile hulls (Theorem~\ref{thm:minkowski}). The gap between them is not an
artefact of any particular router: it equals
$\Delta(\lambda)=\min_k\overline{\,\cdot\,}-\overline{\min_k\,\cdot\,}\ge 0$,
an exchange of minimum and average, and it vanishes exactly when every tile
prefers the same exit (Theorem~\ref{thm:gain}). $\Delta$ is therefore the
\emph{value of content adaptivity}, measurable without building a router.
Experiments on a 42M-parameter neural image codec confirm the characterisation
and show $\Delta/J$ rising from $0.68\%$ to $4.52\%$ as the rate increases,
i.e.\ adaptivity is worth most where the decoder is most expensive.
\end{abstract}

\section{Setup}

\begin{definition}[Multi-exit decode]\label{def:setup}
A frame is partitioned into $N$ disjoint tiles of equal area. A decoder exposes
exits $k\in\mathcal{K}=\{j,\dots,K-1\}$; decoding a tile at exit $k$ costs
$C_k>0$ relative to a full decode and yields squared error $D_i(k)\ge 0$ on
tile~$i$. An \emph{assignment} is a map $a:[N]\to\mathcal{K}$.
\end{definition}

\begin{assumption}[Cost additivity]\label{as:cost}
The compute spent on a frame is $C(a)=\frac1N\sum_{i} C_{a(i)}$.
\end{assumption}

\begin{assumption}[Distortion is a tile mean]\label{as:dist}
The frame distortion is $D(a)=\frac1N\sum_i D_i(a(i))$.
\end{assumption}

Assumption~\ref{as:cost} holds whenever tiles are decoded independently, which
is what a tiled decoder does; Assumption~\ref{as:dist} holds for any per-pixel
loss averaged over the frame, mean squared error included. Both are exact for
the architecture we measure, not approximations. Neither assumes anything about
the $D_i(\cdot)$ themselves --- in particular they need not be convex, monotone,
or identically distributed across tiles.

\section{The achievable set}

Write $\bar D_k=\frac1N\sum_i D_i(k)$ for the mean distortion of exit $k$.

\begin{proposition}[Fixed-proportion mixing]\label{prop:hull}
Let tiles be assigned to exits in proportions $p\in\Delta^{K}$ independently of
their content. The achievable set of $(\Ex C,\Ex D)$ is exactly
$\mathrm{conv}\{(C_k,\bar D_k):k\in\mathcal{K}\}$.
\end{proposition}

\begin{theorem}[Per-tile assignment]\label{thm:minkowski}
The closure of the convex hull of achievable $(C(a),D(a))$ over all assignments
$a$ is the Minkowski average
$\frac1N\bigoplus_{i=1}^{N}\mathrm{conv}\{(C_k,D_i(k)):k\in\mathcal{K}\}$.
\end{theorem}

Theorem~\ref{thm:minkowski} contains Proposition~\ref{prop:hull} as the special
case $D_i\equiv\bar D$. The two sets coincide only in that case; otherwise the
per-tile set is strictly larger, and the next result says by exactly how much.

\section{The value of adaptivity}

For $\lambda\ge 0$ define the Lagrangian values
\begin{equation}\label{eq:J}
J_{\mathrm{ad}}(\lambda)=\frac1N\sum_{i=1}^{N}\min_{k\in\mathcal{K}}
\bigl[D_i(k)+\lambda C_k\bigr],
\qquad
J_{\mathrm{fix}}(\lambda)=\min_{k\in\mathcal{K}}
\bigl[\bar D_k+\lambda C_k\bigr].
\end{equation}

\begin{theorem}[Adaptivity gain]\label{thm:gain}
Under Assumptions~\ref{as:cost}--\ref{as:dist}, for every $\lambda\ge0$,
\begin{equation}\label{eq:gain}
\Delta(\lambda)\;:=\;J_{\mathrm{fix}}(\lambda)-J_{\mathrm{ad}}(\lambda)\;\ge\;0,
\end{equation}
with equality if and only if some $k^\star\in\arg\min_k[D_i(k)+\lambda C_k]$ is
common to all $i\in[N]$.
\end{theorem}

$\Delta(\lambda)$ is computable from the $D_i(k)$ alone. It upper-bounds what
\emph{any} router can gain over fixed proportions at that operating point, and
is attained by the oracle assignment $a^\star_\lambda(i)\in\arg\min_k[D_i(k)+\lambda C_k]$.
It therefore separates two questions that are routinely conflated: whether
content adaptivity is worth anything here, and whether a particular router
realises it.

\begin{proposition}[Frontier recovery]\label{prop:legendre}
The lower boundary of either achievable set is recovered from its $J$ by
$D(C)=\sup_{\lambda\ge0}\{J(\lambda)-\lambda C\}$, and a supporting line at an
optimal point has slope $-\lambda$.
\end{proposition}

\begin{corollary}[Sweeping $\lambda$ is exhaustive]\label{cor:sweep}
Every point on the lower boundary is $a^\star_\lambda$ for some $\lambda\ge0$;
no assignment outside this family is efficient.
\end{corollary}

\begin{remark}[What is \emph{not} claimed]\label{rem:notclaimed}
$D$ is convex in $C$ by Theorem~\ref{thm:minkowski}, but the frontier is usually
plotted as decibels against percentage saving, i.e.\ $\mathrm{dB}(S)=10\log_{10}D$
with $S=1-C/C_{K-1}$. Since $\log$ is concave and increasing, $\log\circ D$ need
be neither convex nor concave, and no such claim follows from the above. We find
it convex empirically (Table~\ref{tab:marginal}); we do not prove it.
\end{remark}

\section{Experiments}

We instantiate Definition~\ref{def:setup} on the intra decoder of a
42.2M-parameter learned image codec (DCVC-UF), split into $K=6$ exits over 12
residual blocks with $j=2$, on $128\times128$ tiles, evaluated on 20 CTC frames
(UVG, HEVC class~E). Distortion is measured against the \emph{released} decoder
from the identical bitstream, so the encoder, hyperprior and entropy model are
byte-identical and the only difference is synthesis. Costs $C_k$ come from a
MAC audit and agree with wall-clock to within $1\%$.

\begin{table}[h]\centering\small
\caption{Per-exit distortion, in dB below the released decoder. The penalty for
leaving early grows $3.4\times$ from the lowest to the highest rate.}
\label{tab:exits}
\begin{tabular}{lrrrrr}
\toprule
QP & exit 2 & exit 3 & exit 4 & exit 5 & bpp \\
\midrule
0  & 0.3447 & 0.1516 & 0.0634 & 0.0228 & 0.1936 \\
32 & 0.6152 & 0.3323 & 0.1096 & 0.0465 & 0.2565 \\
63 & 1.1574 & 0.6310 & 0.1581 & 0.0602 & 0.4559 \\
\bottomrule
\end{tabular}
\end{table}

\begin{table}[h]\centering\small
\caption{$\Delta(\lambda)/J_{\mathrm{ad}}(\lambda)$, the adaptivity gain of
Theorem~\ref{thm:gain} in scale-free form. At $\lambda=3\!\times\!10^{-5}$,
where routing is active, it rises monotonically with rate. At
$\lambda=3\!\times\!10^{-4}$ the oracle itself has collapsed onto one exit and
$\Delta$ vanishes, exactly as the equality condition predicts.}
\label{tab:delta}
\begin{tabular}{lrrrr}
\toprule
QP & $\lambda{=}10^{-5}$ & $3{\times}10^{-5}$ & $10^{-4}$ & $3{\times}10^{-4}$\\
\midrule
0  & 0.24\% & 0.68\% & 1.25\% & 0.09\% \\
16 & 0.41\% & 1.16\% & 0.87\% & 0.01\% \\
32 & 0.71\% & 2.08\% & 0.47\% & 0.00\% \\
48 & 0.88\% & 3.52\% & 0.36\% & 0.00\% \\
63 & 0.93\% & \textbf{4.52\%} & 0.39\% & 0.00\% \\
\bottomrule
\end{tabular}
\end{table}

\begin{table}[h]\centering\small
\caption{Marginal dB per additional 5 points of saving, by linear interpolation
along the measured frontier. Increasing everywhere: each further point of
compute saved costs more than the last. This is the empirical convexity of
Remark~\ref{rem:notclaimed}, not a consequence of the theorems.}
\label{tab:marginal}
\begin{tabular}{lrrrr}
\toprule
QP & $10\!\to\!15$ & $15\!\to\!20$ & $20\!\to\!25$ & $25\!\to\!30$ \\
\midrule
0  & 0.0104 & 0.0164 & 0.0249 & 0.0366 \\
32 & 0.0180 & 0.0292 & 0.0406 & 0.0569 \\
63 & 0.0319 & 0.0505 & 0.0799 & 0.1127 \\
\bottomrule
\end{tabular}
\end{table}

Table~\ref{tab:delta} is the paper's empirical point. The value of content
adaptivity is not constant: it is smallest exactly where a uniformly shallower
decoder would already do (low rate, Table~\ref{tab:exits} left column) and
largest where the deep blocks are doing real work (high rate). A method
evaluated only at low rate would understate its own contribution by a factor of
six.

\section{Proofs}

\begin{proof}[Proof of Proposition~\ref{prop:hull}]
Under fixed proportions $p$, tile assignment is independent of $i$, so
$\Ex C=\sum_k p_kC_k$ and, by Assumption~\ref{as:dist},
$\Ex D=\sum_k p_k\bar D_k$. The map $p\mapsto(\Ex C,\Ex D)$ is affine, hence
maps the simplex $\Delta^K$ onto the convex hull of the images of its vertices,
which are the points $(C_k,\bar D_k)$.
\end{proof}

\begin{proof}[Proof of Theorem~\ref{thm:minkowski}]
By Assumptions~\ref{as:cost}--\ref{as:dist} the achievable pairs are
$\bigl\{\frac1N\sum_i v_i : v_i\in S_i\bigr\}$ with
$S_i=\{(C_k,D_i(k)):k\in\mathcal{K}\}$, i.e.\ the Minkowski average
$\frac1N\bigoplus_i S_i$. Taking convex hulls commutes with Minkowski sums,
$\mathrm{conv}(A\oplus B)=\mathrm{conv}(A)\oplus\mathrm{conv}(B)$, and with
scaling; applying this $N-1$ times gives the claim.
\end{proof}

\begin{proof}[Proof of Theorem~\ref{thm:gain}]
Fix $\lambda\ge0$. For every $k\in\mathcal{K}$,
\[
\bar D_k+\lambda C_k
=\frac1N\sum_{i=1}^{N}\bigl[D_i(k)+\lambda C_k\bigr]
\;\ge\;\frac1N\sum_{i=1}^{N}\min_{k'\in\mathcal{K}}
\bigl[D_i(k')+\lambda C_{k'}\bigr]
=J_{\mathrm{ad}}(\lambda),
\]
since each summand dominates the corresponding minimum. The right-hand side
does not depend on $k$, so taking the minimum over $k$ on the left preserves the
inequality and yields $J_{\mathrm{fix}}(\lambda)\ge J_{\mathrm{ad}}(\lambda)$,
which is \eqref{eq:gain}.

For equality: if some $k^\star$ attains the per-tile minimum for every $i$, then
choosing $k=k^\star$ makes every summand equal to its minimum and the inequality
is tight. Conversely, suppose equality holds and let $k^\star$ attain
$J_{\mathrm{fix}}$. Then
$\frac1N\sum_i\bigl([D_i(k^\star)+\lambda C_{k^\star}]-\min_{k'}[D_i(k')+\lambda C_{k'}]\bigr)=0$.
Every summand is non-negative, so every summand is zero, i.e.\ $k^\star$ attains
the minimum for each $i$.
\end{proof}

\begin{proof}[Proof of Proposition~\ref{prop:legendre} and
Corollary~\ref{cor:sweep}]
Both achievable sets are, after taking convex hulls, compact and convex
(Proposition~\ref{prop:hull}, Theorem~\ref{thm:minkowski}). For such a set the
function $C\mapsto\min\{D:(C,D)\text{ achievable}\}$ is convex and lower
semicontinuous, and $J(\lambda)=\min_{(C,D)}\{D+\lambda C\}$ is its concave
conjugate; the stated inversion is the Fenchel biconjugate applied to a closed
convex function. A minimiser of $D+\lambda C$ lies where a line of slope
$-\lambda$ supports the set, giving the slope claim. Corollary~\ref{cor:sweep}
follows because every exposed point of a compact convex set in $\R^2$ is the
unique minimiser of some linear functional, and non-exposed boundary points are
limits of exposed ones.
\end{proof}

\begin{remark}[Scope]
Theorem~\ref{thm:gain} bounds what adaptivity is worth, not what is realisable
by a decoder that must \emph{infer} the assignment. In our system that
distinction is removed by signalling $a^\star_\lambda$ in the bitstream at
$\lceil\log_2|\mathcal{K}|\rceil$ bits per tile, entropy-coded to
$1.3\times10^{-4}$~bpp --- four orders of magnitude below the rates in
Table~\ref{tab:exits}.
\end{remark}

\end{document}
